\documentclass[11pt]{article}

\usepackage[margin=1in]{geometry}
\usepackage{amsmath,amssymb}
\usepackage[T1]{fontenc}
\usepackage{lmodern}
\usepackage{hyperref}
\usepackage{array}
\usepackage{booktabs}

\title{Technical Description of the Asymmetric Clausius Constraint Analysis}
\author{}
\date{\today}

\begin{document}
\maketitle

\begin{abstract}
This document explains, in plain technical English, what the asymmetric
Clausius constraint analysis does in the
\texttt{nuclear\_matter\_workflows} repository, how the physics is encoded in
the code, which empirical constraints are imposed, and how those constraints
are introduced numerically. The main target is the zero-temperature fitting
workflow in
\texttt{asymmetric\_clausius\_fit/src/asymmetric\_clausius\_fit/solver.py},
which determines the asymmetric Clausius interaction parameters
$a_n$, $a_{pn}$, $b_n$, $b_{pn}$, and $c$ from saturation and symmetry-energy
constraints. The document also explains how the resulting fit table is used as
input for the fixed-$y$ quantum critical-point solver and for the quarkyonic
equation-of-state workflow.
\end{abstract}

\section{Purpose of the Code}

The asymmetric Clausius constraint analysis is the upstream parameter-fitting
stage for several later workflows in this repository.

Its purpose is to construct a family of asymmetric nuclear-matter equations of
state that:
\begin{itemize}
\item reproduce standard symmetric-matter saturation properties;
\item span an empirical window of incompressibility values $K_0$;
\item reproduce the symmetry energy $J$ at saturation;
\item reproduce the symmetry-energy slope $L$ at saturation.
\end{itemize}

The output of the workflow is therefore not a critical point by itself. Rather,
it is a table of interaction parameters for asymmetric Clausius matter, one row
for each target value of $K_0$. Those fitted parameters are then used by the
finite-temperature fixed-$y$ solver and by the quarkyonic EOS workflow.

\section{Relevant Files}

The main files involved in this workflow are listed in
Table~\ref{tab:files}.

\begin{table}[h]
\centering
\begin{tabular}{>{\raggedright\arraybackslash}p{0.42\textwidth} >{\raggedright\arraybackslash}p{0.48\textwidth}}
\toprule
File & Role \\
\midrule
\texttt{notebooks/asymmetric\_clausius\_constraint\_analysis.ipynb} & Main notebook for the zero-temperature constraint analysis, tables, and plots. \\
\texttt{asymmetric\_clausius\_fit/src/asymmetric\_clausius\_fit/constants.py} & Physical constants, empirical targets, and numerical scan settings. \\
\texttt{asymmetric\_clausius\_fit/src/asymmetric\_clausius\_fit/fermi\_gas.py} & Relativistic ideal Fermi-gas relations at $T=0$. \\
\texttt{asymmetric\_clausius\_fit/src/asymmetric\_clausius\_fit/solver.py} & Core solver for symmetric Clausius matter, asymmetric splitting, and constraint enforcement. \\
\texttt{asymmetric\_clausius\_fit/src/asymmetric\_clausius\_fit/finite\_temperature.py} & Downstream fixed-$y$ quantum critical-point workflow that consumes the fitted parameter table. \\
\texttt{asymmetric\_clausius\_fit/src/asymmetric\_clausius\_fit/quarkyonic.py} & Downstream quarkyonic EOS workflow that also consumes the fitted parameter table. \\
\texttt{results/asymmetric\_clausius\_fit\_results.csv} & Main output table of fitted asymmetric Clausius parameters and constraint residuals. \\
\bottomrule
\end{tabular}
\caption{Core files used by the asymmetric Clausius constraint analysis.}
\label{tab:files}
\end{table}

\section{High-Level Workflow}

At a high level, the analysis proceeds as follows:

\begin{enumerate}
\item Choose a target incompressibility $K_0$ in the empirical window.
\item Solve the \emph{symmetric} Clausius model at saturation and determine the
common parameter $c$ that reproduces the chosen $K_0$.
\item At that same symmetric point, recover the corresponding average
parameters $a_{\mathrm{avg}}$ and $b_{\mathrm{avg}}$.
\item Split those averages into asymmetric interaction parameters
$a_n$, $a_{pn}$, $b_n$, and $b_{pn}$.
\item Adjust the attractive split so that the symmetry energy satisfies
$J = S(n_0)$.
\item Adjust the excluded-volume split so that the symmetry slope satisfies
$L = 3 n_0 \, dS/dn|_{n_0}$, while re-solving the attractive split to keep
$J$ fixed.
\item Save the resulting parameter set and diagnostic residuals.
\end{enumerate}

This means that the workflow has a clear nested structure:
\begin{itemize}
\item a symmetric saturation solve;
\item a scan in the common Clausius parameter $c$ to match $K_0$;
\item an asymmetric splitting solve to match $J$ and $L$.
\end{itemize}

\section{Physics Model}

\subsection{Symmetric Clausius Matter}

For symmetric Clausius matter, the code uses the same basic structure as the
existing ground-state workflow:
\begin{align}
f(n) &= 1 - b n, \\
n_{\mathrm{id}} &= \frac{n}{f(n)}, \\
\epsilon(n) &= f(n)\,\epsilon_{\mathrm{id}}(n_{\mathrm{id}})
- \frac{a n^2}{1 + c n}, \\
P(n) &= p_{\mathrm{id}}(n_{\mathrm{id}})
- \frac{a n^2}{(1 + c n)^2}.
\end{align}

Here:
\begin{itemize}
\item $b$ controls the excluded-volume correction;
\item $a$ controls the strength of the attraction;
\item $c$ modifies the density dependence of the attractive Clausius term.
\end{itemize}

The ideal-gas energy density and pressure are evaluated with relativistic
zero-temperature Fermi-gas expressions.

\subsection{Asymmetric Clausius Matter}

For asymmetric matter, the total baryon density is split into proton and
neutron sectors:
\begin{align}
n &= n_p + n_n, \\
n_p &= y n, \\
n_n &= (1-y)n,
\end{align}
where $y$ is the proton fraction.

The excluded-volume structure is made species-dependent:
\begin{align}
x_p &= b_n n_p + b_{pn} n_n, \\
x_n &= b_{pn} n_p + b_n n_n,
\end{align}
with free-volume factors
\begin{align}
f_p &= 1 - x_p, \\
f_n &= 1 - x_n.
\end{align}

If either $f_p \le 0$ or $f_n \le 0$, the state is rejected as unphysical.

The asymmetric Clausius energy density is then
\begin{equation}
\epsilon(n,y)
= f_p \, \epsilon^{\mathrm{id}}_p\!\left(\frac{n_p}{f_p}\right)
+ f_n \, \epsilon^{\mathrm{id}}_n\!\left(\frac{n_n}{f_n}\right)
- \frac{a_n(n_p^2+n_n^2) + 2 a_{pn} n_p n_n}{1 + c n}.
\end{equation}

This means that the code distinguishes:
\begin{itemize}
\item same-species couplings through $a_n$ and $b_n$;
\item proton-neutron couplings through $a_{pn}$ and $b_{pn}$;
\item a common Clausius denominator parameter $c$ shared by all channels.
\end{itemize}

\section{Empirical Constraints}

The entire purpose of the analysis is to impose a specific set of physical
constraints.

\subsection{Symmetric Saturation Constraints}

The symmetric sector is forced to reproduce the standard saturation point:
\begin{align}
n_0 &= 0.16 \ \mathrm{fm}^{-3}, \\
\left.\frac{E}{A} - m\right|_{n_0} &= -16 \ \mathrm{MeV}, \\
P(n_0) &= 0.
\end{align}

These conditions determine the symmetric Clausius parameters at fixed $c$.
Operationally:
\begin{itemize}
\item the zero-pressure condition is used to express $a$ in terms of $b$ and
$c$;
\item the binding-energy condition is then used to solve for $b$.
\end{itemize}

\subsection{Incompressibility Constraint}

The incompressibility is defined as
\begin{equation}
K_0 = 9 \left.\frac{dP}{dn}\right|_{n_0}.
\end{equation}

In this workflow, $K_0$ is the control parameter. The code scans over the
common Clausius parameter $c$ and keeps the value that makes the symmetric
solution reproduce the chosen target $K_0$.

The default target set is
\[
K_0 = 250,\ 260,\ 270,\ 280,\ 290,\ 300,\ 310,\ 315 \ \mathrm{MeV}.
\]

\subsection{Symmetry-Energy Constraint}

The code adopts the parabolic symmetry-energy convention
\begin{equation}
S(n) = E_{\mathrm{PNM}}(n) - E_{\mathrm{SNM}}(n),
\end{equation}
where:
\begin{itemize}
\item PNM means pure neutron matter, corresponding to $y=0$;
\item SNM means symmetric nuclear matter, corresponding to $y=0.5$.
\end{itemize}

The symmetry energy at saturation is then
\begin{equation}
J = S(n_0).
\end{equation}

The target value used by the code is
\[
J = 32.5 \ \mathrm{MeV}.
\]

\subsection{Symmetry-Slope Constraint}

The symmetry-energy slope is defined as
\begin{equation}
L = 3 n_0 \left.\frac{dS}{dn}\right|_{n_0}.
\end{equation}

The target value used by the code is
\[
L = 58.9 \ \mathrm{MeV}.
\]

\section{How the Constraints Are Introduced}

This is the most important conceptual part of the workflow.

\subsection{Step 1: Solve the Symmetric Problem at Fixed \texorpdfstring{$c$}{c}}

For a chosen $c$, the code:
\begin{enumerate}
\item computes the symmetric ideal density at $n_0$;
\item uses the zero-pressure condition to recover $a(b,c)$;
\item solves the binding-energy condition to determine $b$;
\item computes the resulting symmetric point and its $K_0$ value.
\end{enumerate}

If the excluded-volume fraction becomes non-positive, or if the scan in $b$
does not produce a valid sign change for the binding residual, that trial point
is discarded.

\subsection{Step 2: Use \texorpdfstring{$K_0$}{K0} to Fix the Common Clausius Parameter}

Once the symmetric point exists for a given $c$, the code computes
\[
K_0(c) = 9 \left.\frac{dP}{dn}\right|_{n_0}.
\]

The solver then scans the interval
\[
0 \le c \le 4.74 \ \mathrm{fm}^3
\]
and looks for a sign change in
\[
K_0(c) - K_0^{\mathrm{target}}.
\]

When such a bracket is found, a bisection solve returns the value of $c$ that
reproduces the requested incompressibility.

\subsection{Step 3: Introduce Asymmetry Through Average-and-Split Variables}

After the symmetric point has been determined, the asymmetric parameters are
constructed from average and split variables:
\begin{align}
a_n &= a_{\mathrm{avg}} - \Delta a, \\
a_{pn} &= a_{\mathrm{avg}} + \Delta a, \\
b_n &= b_{\mathrm{avg}} - \Delta b, \\
b_{pn} &= b_{\mathrm{avg}} + \Delta b.
\end{align}

This is a deliberate modeling choice:
\begin{itemize}
\item the symmetric saturation fit fixes the averages;
\item the isovector information is carried by the splits $\Delta a$ and
$\Delta b$;
\item the Clausius parameter $c$ is left common to all channels.
\end{itemize}

\subsection{Step 4: Solve \texorpdfstring{$\Delta a$}{Delta a} from the \texorpdfstring{$J$}{J} Constraint}

For a given trial $\Delta b$, the code scans $\Delta a$ and evaluates
\[
J(\Delta a,\Delta b) - J_{\mathrm{target}}.
\]

This means that the attractive split is used to tune the symmetry energy at
saturation. If the resulting parameters make any of
$a_n$, $a_{pn}$, $b_n$, or $b_{pn}$ non-positive, the trial point is discarded.

\subsection{Step 5: Solve \texorpdfstring{$\Delta b$}{Delta b} from the \texorpdfstring{$L$}{L} Constraint}

The code then scans $\Delta b$. For each trial $\Delta b$:
\begin{enumerate}
\item it first re-solves $\Delta a$ so that $J$ remains fixed;
\item it computes the resulting $L$ value;
\item it looks for a sign change in
\[
L(\Delta a(\Delta b),\Delta b) - L_{\mathrm{target}}.
\]
\end{enumerate}

Once a bracket is found, a bisection solve returns the final $\Delta b$, and
then $\Delta a$ is solved once more at that final value.

So the hierarchy of constraints is:
\begin{itemize}
\item $K_0$ fixes the common Clausius parameter $c$;
\item $J$ fixes the attractive split $\Delta a$;
\item $L$ fixes the excluded-volume split $\Delta b$, with $J$ enforced at each
step.
\end{itemize}

\section{Numerical Structure of the Solver}

The code uses simple local numerical routines rather than external libraries.

\subsection{Bracketing and Bisection}

All major constraints are imposed through sign-change scans followed by
bisection:
\begin{itemize}
\item the symmetric binding residual is bracketed in $b$;
\item the incompressibility residual is bracketed in $c$;
\item the $J$ residual is bracketed in $\Delta a$;
\item the $L$ residual is bracketed in $\Delta b$.
\end{itemize}

This makes the algorithm robust and easy to interpret physically.

\subsection{Numerical Derivatives}

Two derivative-based quantities are evaluated numerically:
\begin{align}
K_0 &= 9 \left.\frac{dP}{dn}\right|_{n_0}, \\
L &= 3 n_0 \left.\frac{dS}{dn}\right|_{n_0}.
\end{align}

The derivatives are computed with local finite differences using step sizes
stored in the constants file.

\subsection{Saturation Check}

After a symmetric solution has been built, the code performs an additional
consistency check by minimizing
\[
\frac{E}{A} - m
\]
in a window around saturation. This returns:
\begin{itemize}
\item the numerical saturation density $n_{\mathrm{sat}}$;
\item the corresponding energy per particle;
\item the pressure at the minimum.
\end{itemize}

These values are stored in the output table as diagnostics.

\section{Control Flow of the Main Public Functions}

\subsection{\texttt{symmetric\_clausius\_point\_from\_c}}

This function solves the symmetric Clausius model at fixed $c$ and returns the
corresponding symmetric point. It is the building block for all later stages.

\subsection{\texttt{find\_c\_for\_target\_k0}}

This function scans the Clausius parameter $c$ and finds the value that
reproduces the requested incompressibility.

\subsection{\texttt{fit\_asymmetric\_clausius\_target}}

This is the main single-case solver. It:
\begin{enumerate}
\item fixes $c$ from the target $K_0$;
\item reconstructs the symmetric point and its average parameters;
\item solves the asymmetric splits from $J$ and $L$;
\item packages the full result, including residuals and diagnostics.
\end{enumerate}

This is the function that a user should think of as
``fit one asymmetric Clausius EOS for one target incompressibility''.

\subsection{\texttt{compute\_target\_family}}

This function simply repeats the single-case fit across the full list of target
$K_0$ values and returns a family of fitted asymmetric Clausius models.

\section{What the Output Table Contains}

The main output CSV stores, for each target $K_0$:
\begin{itemize}
\item the common Clausius parameter $c$;
\item the symmetric averages $a_{\mathrm{avg}}$ and $b_{\mathrm{avg}}$;
\item the asymmetric parameters $a_n$, $a_{pn}$, $b_n$, and $b_{pn}$;
\item the ratios $a_{pn}/a_n$ and $b_{pn}/b_n$;
\item the recovered values of $J$, $L$, and $K_0$;
\item the residuals relative to the target constraints;
\item saturation diagnostics such as $n_{\mathrm{sat}}$ and $p_{\mathrm{sat}}$.
\end{itemize}

This makes it easy to check both the physics and the numerical quality of the
fit.

\section{How the Fitted Parameters Are Used Later}

The asymmetric Clausius constraint analysis is not an isolated calculation. It
feeds directly into later physics workflows.

\subsection{Fixed-\texorpdfstring{$y$}{y} Quantum Critical Points}

The file \texttt{finite\_temperature.py} reads one row of the fit table and
uses the fitted parameters
\[
(a_n,\ a_{pn},\ b_n,\ b_{pn},\ c)
\]
to build the finite-temperature pressure at fixed proton fraction $y$.

That downstream solver then computes the quantum critical point
\[
(T_c,\ n_c,\ P_c)
\]
from the conditions
\[
\left(\frac{\partial P}{\partial n}\right)_T = 0,
\qquad
\left(\frac{\partial^2 P}{\partial n^2}\right)_T = 0.
\]

So the asymmetric constraint analysis provides the interaction input for the
SCF solver described in the separate fixed-$y$ quantum note.

\subsection{Quarkyonic EOS Construction}

The file \texttt{quarkyonic.py} also consumes the fitted table. It converts the
parameters from fm/MeV units to GeV-based units and uses them to construct the
hadronic part of the quarkyonic EOS.

In that sense, the asymmetric Clausius fit is the hadronic calibration stage
for the quarkyonic calculation.

\section{Why the Code Rejects Some Trial Points}

The solver returns \texttt{None} when a trial state is unphysical or numerically
unsafe. Typical failure modes are:
\begin{itemize}
\item the excluded-volume fraction becomes non-positive;
\item the Clausius denominator $1 + c n$ becomes non-positive;
\item no sign-change bracket is found for one of the required residuals;
\item one of the split parameters becomes non-positive.
\end{itemize}

This makes the workflow conservative. It prefers discarding an invalid trial
point to forcing a parameter set with poor physical meaning.

\section{Summary}

The asymmetric Clausius constraint analysis is the zero-temperature fitting
stage that calibrates asymmetric nuclear-matter interactions before any
finite-temperature or quarkyonic calculation is performed.

Its essential logic is:
\begin{enumerate}
\item solve symmetric Clausius matter at saturation;
\item use the common Clausius parameter $c$ to match the target
incompressibility $K_0$;
\item split the symmetric averages into same-species and proton-neutron
interaction parameters;
\item tune the attractive split to reproduce $J$;
\item tune the excluded-volume split to reproduce $L$ while keeping $J$ fixed;
\item store the resulting parameter family for downstream workflows.
\end{enumerate}

Therefore, the answer to
``what are the constraints and how are they introduced?'' is:
\begin{itemize}
\item \textbf{symmetric saturation constraints:}
$P(n_0)=0$ and $(E/A-m)|_{n_0}=-16\ \mathrm{MeV}$ fix the symmetric point at
fixed $c$;
\item \textbf{incompressibility constraint:}
$K_0$ fixes the common Clausius parameter $c$;
\item \textbf{symmetry-energy constraint:}
$J=S(n_0)$ fixes the attractive split $\Delta a$;
\item \textbf{symmetry-slope constraint:}
$L=3n_0\,dS/dn|_{n_0}$ fixes the excluded-volume split $\Delta b$.
\end{itemize}

These constraints are the conceptual heart of the asymmetric Clausius fitting
workflow, and the resulting parameter table is the bridge to both the
finite-temperature SCF critical-point analysis and the quarkyonic EOS
construction.

\end{document}
